% Imporditud: tools/import_eksperimendid.py
% Allikas: Eesti füüsikaolümpiaadi eksperimendi ülesannete
%   kogu (Rael Kalda, 2023), ülesanne Ü30, lahendus L30.
\setAuthor{Erkki Tempel}
\setRound{piirkonnavoor}
\setYear{2022}
\setNumber{P E1}
\setDifficulty{1}
\setTopic{staatika}
\setVahendid{joonlaud, A4 paber (80 g/m$^2$)}
\setPunktid{10}
\setAllikas{Eksperimendi kogu Ü30, lahendus lk 46}
\setLahendusseis{toores}

\prob{Joonlaua mass}
Määrake joonlaua mass.

\hint

\solu
Mõõdame A4 paberi pikkuse ja laiuse ning teades paberi pindtihedust, saame leida

A4 paberi massi mp $\approx$ 5 g

Leiame joonlaua masskeskme laua serval.

Asetame kokkuvolditud A4 paberi joonlaua ühe otsa peale ning leiame joonlaua ja

paberi tasakaalupunkti.

Mõõdame paberi keskkoha kauguse toetuspunktist lp ning joonlaua masskeskme

kauguse toetuspunktist lj.

Kasutades kangireeglit, saame tasakaalu jaoks kirjutada seose ning leida joonlaua

massi mj.

mplp = mj lj $\Rightarrow$ mj = mplp lj

Hindamine: Paberi massi leidmine --- 2 p. Joonlaua masskeskme leidmine --- 2 p. Joonlaua tasakaalupunkti leidmine koos paberiga --- 1 p. Tasakaalu seoste avaldamine --- 1 p. Joonlaua massi avaldamine --- 1 p. Jõuõlgade lp ja lj korrektne mõõtmine koos kordusmõõtmistega --- 2 p. Joonlaua massi arvutamine --- 1 p.
\probend
